% =====================================================================
% Section 6 --- Beyond demand-only shocks
% ---------------------------------------------------------------------
% Job: what simplifies and where variety lives --- the equivalences and
%      their scope (demand-only, the A = 1 branch the numeraire selects);
%      the sectoral family and its exact nesting (the rigid corner IS the
%      eta = 0 endpoint); the executed ladder and rigid-group cells; the
%      supply-shock arm as identification vehicle; the no-go corollary; BF
%      repositioned here as the frictionless pole.
% Mindmap: N03, N13, N02D.
% Status: STUB --- narrative outline v3 (2026-09-20).
% =====================================================================
\section{Beyond demand-only shocks}\label{sec:beyond-demand-only}

% TODO: scope the equivalences airtight; give a clean supply-shock
% counterexample and use it to motivate this section rather than hide the
% boundary. Reference paper/framing_gamma.tex.

\subsection{The equivalences and their scope}\label{subsec:equivalences}
% TODO: Lemma 1 (price-block demand-invariance), Lemma 2 (nonsubstitution,
% Samuelson 1951), Propositions 1-3; financing neutrality; full proofs in
% App. A.

\subsection{The sectoral family}\label{subsec:sectoral-family}
% TODO: L^cm_i = Lbar_i (w_i/Pi)^{eta_s,i}; the rigid corner is exactly the
% eta = 0 endpoint (Prop. 3); the executed ladder (BETA-F2-etas{025,05,1,2} =
% 0.152897 / 0.105666 / 0.065422 / 0.037171) and the rigid-group cells.

\subsection{The supply-shock arm}\label{subsec:supply-arm}
% TABLE (linked, revision/tables/supply_etas_flows.md): the s1 identification
% vehicle; supply_etas_prog_flows.md / supply_etas_unif_flows.md for the other
% two designs (189 cells total).
% TODO: the identification vehicle (supply_etas_s1); 189 cells across three
% designs; separates the ladder rungs that demand-only shocks cannot.

\subsection{The no-go corollary and the BF pole}\label{subsec:no-go}
% TODO: nominal wage rules are vacuous in both gauges --- the rigidity that
% matters must live in sectoral labour markets. BF repositioned as the
% frictionless pole (peer-modesty: late introduction, no vantage status).
% ===== CARRIED OVER FROM AN OLD DRAFT =============================
% [devised for an old draft; not in the validated submission]
% source: revised_manuscript/chapters/ch05_ces_functions.tex (commit 41144dd)
% =================================================================
% With regard to how changes in assumptions on substitution elasticities impact on final outcomes Figure \ref{fig:elasticity-gradient} indicates that real GDP is typically increasing in all three elasticities employed -- the higher the chosen elasticities, the higher are the respective GDP-estimates. This finding reproduces the intuition that easier substitution will lead to higher (aggregate) output and, hence, CES-based estimates are consistently lower than their Cobb-Douglas counterparts.
% 
% Against this backdrop we concede that the introduction of Leontief-like properties in the production function helps bridging paradigmatic cleavages on a level of conceptual foundations -- e.g. with respect to the question how to adequately represent production processes. However, we find that introducing such a conceptual shift in neoclassical IO-models will actually serve to increase the gap between both approaches in terms of predicted outcomes. This result is not too surprising when taking into account that the assumption of a Leontief technology implies less flexibility as compared to a Cobb-Douglas technology. Hence, in the context of assumptions on substitutability neoclassical approaches are more optimistic than the heterodox archetype: viewed in isolation greater substitutability comes with higher prospects for growth in case of demand shocks and, correspondingly, leads to a more positive assessment of fiscal intervention in terms of forced green investment. However, this relatively more optimistic stance does in no way compensate for the fact that the estimated growth effects remain negative, which is effectively driven by assumptions on capacity constraints. However, before turning to this key aspect in section \ref{sec:labor} further below, we first inspect structural differences between model outcomes on a more granular level.


% ===== CARRIED OVER FROM AN OLD DRAFT =============================
% [devised for an old draft; not in the validated submission]
% source: revised_manuscript/chapters/ch06_labor_slack.tex (commit 41144dd)
% =================================================================
% Nonetheless, on a general level this example again illustrates how the introduction of labor slack can lead to more optimistic results in terms of GDP growth. The key requirement is the willingness to assume that such an expansion of the labor force is indeed a plausible reaction -- either due to labor market considerations or as a general response to increased effective demand. While this feature is not truly surprising it shows how modifying dominant convictions on the role and modeling of capacity constraints indeed opens up a shared realm, in which the simplistic juxtaposition of an 'expansionary' Leontief-approach and a 'contractionary' CGE-perspective becomes subject to a more nuanced perspective that allows for pragmatically interpreting labor slack as a moderating variable to explore some shared ground between these two competing visions of the economy.

